\documentclass[11pt,a4paper]{article}
\input{../../shared/preamble}
\SpeedHeader{Geometry}{3}

\begin{document}

\SpeedTitleBlock{Daily Geometry Drill \#3}{Henry Koduthore}
\SpeedMeta{Two-circle systems, radical axes and common chords}{relative positions of two circles, radical axis by subtraction, common chord length, orthogonal circles, tangency to both coordinate axes}
\noindent\textit{New toolkit today: Relative Positions (4/3/2/1 tangents) $\cdot$ Radical Axis by subtraction $\cdot$ Common Chords $\cdot$ Orthogonal Circles ($d^2=r_1^2+r_2^2$).}

\vspace{2pt}
\noindent\textit{Attempt each question before checking answers. Time yourself. No calculators.}

\section*{Section A \quad Rapid Recognition \hfill{\normalfont\small($\leq$15 s each)}}
% ══════════════════════════════════════════════════════════════════════════════

\noindent\textit{Compare the centre distance $d$ with $r_1+r_2$ and $|r_1-r_2|$, or subtract the equations to read the radical axis.}

\begin{enumerate}

\item The circle $\;x^2+y^2-12x-16y+75=0$ is tangent to the circle $\;x^2+y^2=25$. Write down the number of common tangents of the two circles.

\item The circles with centre $(0,0)$ radius $5$ and centre $(13,0)$ radius $3$ are disjoint. Write down the number of common tangents of the two circles.

\item Write down the radical axis of the circles $\;x^2+y^2-6x+2y-11=0$ and $\;x^2+y^2-2x-4y+1=0$.

\item Write down whether the circles $\;x^2+y^2=9$ and $\;x^2+y^2+6x+8y+9=0$ are orthogonal (yes or no).

\item The circles $\;x^2+y^2=25$ and $\;x^2+y^2-10x+5=0$ intersect in two points. Find the length of their common chord.

\item Two circles have radii $3$ and $5$ and their centres are $4$ units apart. Write down the number of common tangents.

\item The circle $\;x^2+y^2-10x-10y+k=0$ is tangent to both coordinate axes. Find $k$.

\item The circle $\;x^2+y^2+4x-2y+C=0$ is orthogonal to the circle $\;x^2+y^2-4x+4y-1=0$. Find $C$.

\item How many common tangents do the circles $\;x^2+y^2-4x-6y+12=0$ and $\;x^2+y^2+2x-10y+25=0$ have?

\item Two circles have radii $7$ and $3$ and exactly $3$ common tangents. Find the distance between their centres.

\end{enumerate}

% ══════════════════════════════════════════════════════════════════════════════
\section*{Section B \quad Manipulation Drills \hfill{\normalfont\small($\leq$50 s each)}}
% ══════════════════════════════════════════════════════════════════════════════

\noindent\textit{Complete the square, locate the centres, then classify the two-circle interaction or compute the radical axis.}

\begin{enumerate}

\item How many common tangents do the circles $\;x^2+y^2-2x-4y-20=0$ and $\;x^2+y^2-6x-2y+6=0$ have?
\begin{itemize}
\item[A)] $0$
\item[B)] $1$
\item[C)] $2$
\item[D)] $3$
\item[E)] $4$
\end{itemize}

\item The circle $\;x^2+y^2+ax+by+c=0$ passes through the origin, and its radical axis with the circle $\;x^2+y^2+6x-4y+10=0$ is the line $\;2x-y+5=0$. What is $a+b$?
\begin{itemize}
\item[A)] $3$
\item[B)] $5$
\item[C)] $-3$
\item[D)] $0$
\end{itemize}

\item The circle $\;x^2+y^2-4x+2y-5=0$ is orthogonal to $\;x^2+y^2+6x-2y+k=0$. Find $k$.
\begin{itemize}
\item[A)] $-9$
\item[B)] $9$
\item[C)] $-5$
\item[D)] $14$
\end{itemize}

\item A circle is tangent to both coordinate axes in the first quadrant and passes through the point $(1,8)$. Find the smaller possible radius.

\item Two circles of equal radius $5$ have centres $(0,0)$ and $(6,0)$. Their radical axis has equation?
\begin{itemize}
\item[A)] $x=3$
\item[B)] $x=6$
\item[C)] $y=3$
\item[D)] $y=0$
\end{itemize}

\item Two circles of radii $4$ and $6$ have centres $9$ units apart. How many common tangents do they have?
\begin{itemize}
\item[A)] $0$
\item[B)] $1$
\item[C)] $2$
\item[D)] $3$
\item[E)] $4$
\end{itemize}

\item A circle with centre $(7,1)$ cuts the circle $\;x^2+y^2-2x+4y-4=0$ at right angles. What is its radius?
\begin{itemize}
\item[A)] $3\sqrt5$
\item[B)] $3\sqrt5-3$
\item[C)] $3$
\item[D)] $6$
\end{itemize}

\item The circles $\;x^2+y^2=25$ and $\;x^2+y^2-12x-6y+25=0$ meet at $P$ and $Q$. Find the midpoint of $PQ$.

\item A circle with its centre in the first quadrant touches both coordinate axes and the line $\;3x+4y=12$. What is the sum of the possible radii?
\begin{itemize}
\item[A)] $1$
\item[B)] $6$
\item[C)] $7$
\item[D)] $5$
\end{itemize}

\item The point $P=(1,t)$ lies on the radical axis $\;3x-4y+5=0$ of two circles. Find $t$.

\end{enumerate}

% ══════════════════════════════════════════════════════════════════════════════
\section*{Section C \quad Substitution \& Structure \hfill{\normalfont\small($\leq$75 s each)}}
% ══════════════════════════════════════════════════════════════════════════════

\noindent\textit{Hard, non-routine: two or three circle tools at once --- subtract for the radical axis, then turn that line into a chord or a tangency count. No brute intersection.}

\begin{enumerate}

\item Two circles have the same radius. Their centres are $C_1(-2,1)$ and $C_2(3,-2)$. The circles intersect in two distinct points $P$ and $Q$. Which is the equation of the line $PQ$? \textit{\small(after TMUA 2021 Paper 1 Q1)}
\begin{itemize}
\item[A)] $5x+3y=4$
\item[B)] $3x-5y=4$
\item[C)] $5x-3y=4$
\item[D)] $5x-3y=1$
\end{itemize}

\item The circles $C_1:(x-1)^2+y^2=4$ and $C_2:(x-5)^2+(y-3)^2=9$ have radii $2$ and $3$. How many common tangents do $C_1$ and $C_2$ have? \textit{\small(after TMUA 2019 Paper 1 Q7)}
\begin{itemize}
\item[A)] $0$
\item[B)] $1$
\item[C)] $2$
\item[D)] $3$
\item[E)] $4$
\end{itemize}

\item The circles $x^2+y^2=169$ and $x^2+y^2-6x+8y-119=0$ meet at $P$ and $Q$. What is the length $PQ$?
\begin{itemize}
\item[A)] $12$
\item[B)] $16$
\item[C)] $24$
\item[D)] $26$
\end{itemize}

\item A circle in the first quadrant is tangent to both coordinate axes and passes through the point $(9,2)$. Its two possible radii are $r_1<r_2$. What is $r_1+r_2$? \textit{\small(adapted from JZMaths Mock 1 Q5)}
\begin{itemize}
\item[A)] $14$
\item[B)] $17$
\item[C)] $22$
\item[D)] $34$
\end{itemize}

\item The circles $x^2+y^2=25$ and $x^2+y^2-16x-8y-1=0$ meet at $P$ and $Q$. What is the radius of the circle through $P$, $Q$ and the point $(0,1)$?
\begin{itemize}
\item[A)] $9$
\item[B)] $2\sqrt{13}$
\item[C)] $4\sqrt5$
\item[D)] $13$
\end{itemize}

\item A circle has equation $x^2+ax+y^2+by+c=0$ where $a,b,c\neq0$. Which condition is necessary and sufficient for the circle to be tangent to the $y$-axis? \textit{\small(after TMUA 2021 Paper 2 Q15)}
\begin{itemize}
\item[A)] $a^2=4c$
\item[B)] $b^2=4c$
\item[C)] $a^2+b^2=4c$
\item[D)] $c=0$
\end{itemize}

\item The equal circles $C_1:(x+2)^2+(y-1)^2=3$ and $C_2:(x-4)^2+(y-1)^2=3$ have centres $6$ apart. A transverse common tangent with positive gradient makes an acute angle $\theta$ with the $x$-axis. What is $\tan\theta$? \textit{\small(adapted from TMUA 2020 Paper 1 Q16)}
\begin{itemize}
\item[A)] $\dfrac12$
\item[B)] $\dfrac{\sqrt2}{2}$
\item[C)] $\sqrt2$
\item[D)] $\dfrac{\sqrt3}{3}$
\end{itemize}

\item The three circles $C_1:x^2+y^2+4x-1=0$, $C_2:x^2+y^2-8y+11=0$, $C_3:x^2+y^2+2x-12y+13=0$ have pairwise radical axes concurrent at $R$. What are the coordinates of $R$? \textit{\small(adapted from Vantage Mock Paper 2 Q10 $/$ BMO1 2017 Q4)}
\begin{itemize}
\item[A)] $(1,1)$
\item[B)] $(1,-1)$
\item[C)] $(3,1)$
\item[D)] $(1,2)$
\end{itemize}

\end{enumerate}

% ══════════════════════════════════════════════════════════════════════════════
\section*{Section D \quad Challenge Ramp \hfill{\normalfont\small($\leq$90 s each)}}
% ══════════════════════════════════════════════════════════════════════════════

\noindent\textit{Hardest 5: combine radical axes, orthogonality and chord loops end-to-end. The numbers still cancel, but the route is hidden.}

\begin{enumerate}

\item The circles $x^2+y^2-4x-2y-4=0$ and $x^2+y^2-20x-14y+124=0$ have centres $O_1$ and $O_2$. A common external tangent touches them at $P$ and $Q$ respectively. What is the area of the quadrilateral $O_1PQO_2$?
\begin{itemize}
\item[A)] $24$
\item[B)] $16\sqrt6$
\item[C)] $40$
\item[D)] $8\sqrt6$
\end{itemize}

\item The circles $(x+4)^2+(y+1)^2=64$ and $(x-8)^2+(y-4)^2=r^2$ ($r>0$) have exactly one point in common. What is the difference between the two possible values of $r$? \textit{\small(after TMUA 2019 Paper 1 Q6)}
\begin{itemize}
\item[A)] $8$
\item[B)] $16$
\item[C)] $26$
\item[D)] $42$
\end{itemize}

\item Three circles of radii $3$, $4$ and $5$ are pairwise externally tangent. What is the area of the triangle formed by joining their centres? \textit{\small(adapted from BMO1 2017 Q4)}
\begin{itemize}
\item[A)] $12\sqrt5$
\item[B)] $12$
\item[C)] $84$
\item[D)] $30$
\end{itemize}

\item Two circles of radii $3$ and $4$ are orthogonal. What is the length of their common chord?
\begin{itemize}
\item[A)] $\dfrac{12}{5}$
\item[B)] $\dfrac{24}{5}$
\item[C)] $6$
\item[D)] $\dfrac{48}{5}$
\end{itemize}

\item \begin{minipage}[t]{0.60\linewidth}
The circles $x^2+y^2-8y=0$ and $x^2+y^2-24x-18y+144=0$ both touch the $x$-axis and touch each other. A third circle lies in the gap between them and the $x$-axis, and touches both circles and the $x$-axis. What is its radius?
\begin{itemize}[itemsep=5pt]
\item[A)] $\dfrac65$
\item[B)] $\dfrac{36}{13}$
\item[C)] $\dfrac{169}{100}$
\item[D)] $\dfrac{36}{25}$
\end{itemize}
\end{minipage}\hfill
\begin{minipage}[t]{0.34\linewidth}
\vspace{0pt}
\centering
\begin{tikzpicture}[scale=0.18]
\draw[->] (-4.5,0)--(21.5,0) node[right]{\small $x$};
\draw[->] (0,-0.8)--(0,9.5) node[above]{\small $y$};
\draw[speedgeom] (0,4) circle (4);
\draw[speedgeom] (12,9) circle (9);
\filldraw[fill=gray!15, draw=darkgray] (4.8,1.44) circle (1.44);
\end{tikzpicture}
\end{minipage}

\end{enumerate}

\SpeedClosing{``Two circles are decoded by one subtraction: locate each centre, subtract the equations, then compare $d$ with $r_1+r_2$ and $|r_1-r_2|$.''}

\end{document}
